The reconstructed OOF evaluation does not support a uniform advantage of Full over all alternatives (Tables~\ref{tab:whole_volume_generation},\ref{tab:paired_generation}). Full outperformed the clinically conditioned regression U-Net across the four whole-volume metrics for both tracers, while the magnitude of the AV45 MSE effect was small ($|d_z|=0.075$). Relative to TabularMLP, Full had higher SSIM but lower PSNR and higher MSE for both tracers; MAE did not differ after the specified Holm correction. TracerOnly had better FDG results on all four metrics and higher AV45 SSIM than Full. Single-tracer DIReCT had higher FDG SSIM, whereas Full had better FDG PSNR/MSE and AV45 SSIM. These independently trained models differ in initialization, conditioning or implementation, so the comparisons do not isolate a single causal component or prove multi-tracer synergy.

In the masked evaluation (Tables~\ref{tab:r3_11_masked_summary},\ref{tab:r3_11_masked_paired}), Full had higher mean SSIM/PSNR and lower mean MSE/MAE than Clinical U-Net and MRI-only Swin-UNETR for both tracers; all 16 predefined comparisons met $p_{\mathrm{Holm}}<0.05$ within their two separate eight-test families. The masked cohort contained 1,597 FDG and 1,364 AV45 observations after the existing MRI/mask QC exclusion. For Clinical U-Net AV45 MSE/MAE, the standardized effects were small ($|d_z|=0.071/0.076$), and only 56.4--57.8\% of subjects favored Full. Thus statistical support does not imply uniform individual or clinical benefit. Full/Clinical U-Net used epoch 165, while MRI-only Swin-UNETR used epoch 200; information and training endpoints differ.

Separate whole-volume and masked summaries are retained because background exclusion changes the metric definitions and values. Supplemental FDG and AV45 examples use the median, 10th-percentile and worst QC-passed subjects selected by saved whole-volume SSIM, and show MRI, real PET, synthetic PET and absolute-error maps (Figs.~\ref{fig:cases_fdg},\ref{fig:cases_av45}).
